\section{Ordinary Differential Equations}

\subsection{Basics of one-step methods}

\begin{exercise}
	For each IVP, solve the problem using Euler's method. (i) Plot the solution
	for $n=320$. (ii) For $n=10\cdot2^k$, $k=2,3,\ldots,10$, compute the error as
	$||u-\hat{u}||_\infty=\max_{0\leq i\leq n} |u_i-\hat{u}(t_i)|$ and at the final
	time, $|u_n-\hat{u}(t_n)|$. Make a log-log convergence plot, including a
	reference line for first-order convergence.
	\begin{enumerate}
		\item $u' = -2t u$, $\ 0 \le t \le 2$, $\ u(0) = 2$; \\
		      \[
			      \hat{u}(t) = 2e^{-t^2}
		      \]

		\item $u' = u + t$, $\ 0 \le t \le 1$, $\ u(0) = 2$; \\
		      \[
			      \hat{u}(t) = -1-t+3e^t
		      \]

		\item $(1+t^3)uu' = t^2$, $\ 0 \le t \le 3$, $\ u(0) =1$; \\
		      \[
			      \hat{u}(t) =[1+(2/3)\ln (1+t^3)]^{1/2}
		      \]
	\end{enumerate}
\end{exercise}

\begin{solution}

	The explicit Euler method is a first-order Runge-Kutta numerical integration
	scheme for solving initial value problems (IVPs) of the form $u' = f(t, u)$.
	Given a temporal domain $[t_0, t_f]$ discretized into $n$ steps of uniform size
	$h = (t_f - t_0)/n$, the system is iteratively updated via:
	\[
		u_{i+1} = u_i + h f(t_i, u_i)
	\]
	Because the method relies on a truncated Taylor series approximation, the
	local truncation error introduced at each discrete step scales as
	$\mathcal{O}(h^2)$. However, as these errors accumulate sequentially over $n
		\propto h^{-1}$ steps, the global error bounds scale as $\mathcal{O}(h)$,
	mathematically equivalent to $\mathcal{O}(n^{-1})$.

	To empirically verify this theoretical first-order convergence, we compute
	explicit Euler solutions across progressively refined grids where $n = 10 \cdot
		2^k$ for $k \in \{2, 3, \dots, 10\}$. Two distinct error metrics are recorded:
	\begin{enumerate}
		\item The global maximum absolute divergence ($L^\infty$ norm):
		      $||u-\hat{u}||_\infty = \max_{i} |u_i-\hat{u}(t_i)|$.
		\item The final accumulated displacement at the domain boundary:
		      $E_{\text{final}} = |u_n-\hat{u}(t_n)|$.
	\end{enumerate}

	\paragraph{Results \& Discussion}
	For the baseline resolution of $n=320$, the numerical solutions smoothly
	trace the analytical trajectories for all three IVPs. Because $h$ is
	sufficiently small, virtually zero visually distinguishable deviation is
	observed on linear scales.

	Both the maximum absolute error arrays ($L^\infty$) and the final boundary
	error arrays run parallel to the theoretical $\mathcal{O}(n^{-1})$ reference
	lines. By definition, the final-time error is bounded from above by the
	$L^\infty$ norm ($|u_n - \hat{u}(t_n)| \leq ||u-\hat{u}||_\infty$).

	\paragraph{Figures}
	The transient numerical solutions at $n=320$ (left column) and the
	corresponding log-log $L^\infty$ vs. boundary convergence analyses (right
	column) are presented below.

	\vspace{1em}
	\begin{center}
		\begin{tabular}{cc}
			\includegraphics[width=0.45\textwidth]{figures/6-2/euler-method-1.pdf} &
			\includegraphics[width=0.45\textwidth]{figures/6-2/euler-method-error-1.pdf}                                              \\
			{\small (1a) Solution for $u' = -2t u$}                                & {\small (1b) Convergence for $u' = -2tu$}        \\[1.5em]
			\includegraphics[width=0.45\textwidth]{figures/6-2/euler-method-2.pdf} &
			\includegraphics[width=0.45\textwidth]{figures/6-2/euler-method-error-2.pdf}                                              \\
			{\small (2a) Solution for $u' = u + t$}                                & {\small (2b) Convergence for $u' = u + t$}       \\[1.5em]
			\includegraphics[width=0.45\textwidth]{figures/6-2/euler-method-3.pdf} &
			\includegraphics[width=0.45\textwidth]{figures/6-2/euler-method-error-3.pdf}                                              \\
			{\small (3a) Solution for $(1+t^3)uu' = t^2$}                          & {\small (3b) Convergence for $(1+t^3)uu' = t^2$} \\
		\end{tabular}
	\end{center}

\end{solution}

\begin{exercise}
	Solve the following IVPs using Euler's method with $n=1000$ steps. Plot the
	solution and its first derivative together on one plot, and plot the error in
	each component as functions of time on another.
	\begin{enumerate}
		\item $y''+ 9y = \sin(2t)$, $\ 0< t< 2\pi$, $\ y(0) = 2$, $\ y'(0) = 1$; \\
		      \[
			      \hat{y}(t) = (1/5) \sin(3t) + 2 \cos (3t)+  (1/5) \sin (2t)
		      \]
		\item $y''- 4y = 4t$, $\: 0< t< 1.5$, $\: y(0) = 2$, $\: y'(0) = -1$; \\
		      \[
			      \hat{y}(t) = e^{2t} + e^{-2t}-t
		      \]
		\item $y''+ 4y'+ 4y = t$, $\: 0< t< 4$, $\: y(0) = 1$,$\: y'(0) = 3/4$; \\
		      \[
			      \hat{y}(t) = (3t+5/4)e^{-2t} + (t-1)/4
		      \]
	\end{enumerate}
\end{exercise}

\begin{solution}

	To apply the first-order forward Euler method, each second-order ordinary
	differential equation (ODE) must be reduced to a system of coupled first-order
	ODEs. We introduce the state vector $\mathbf{u}(t) = [u_1(t), u_2(t)]^\intercal
			= [y(t), y'(t)]^\intercal$. The general second-order ODE $y'' = f(t, y, y')$ can
	thus be expressed as:
	\[
		\mathbf{u}'(t) =
		\begin{bmatrix}
			u_1'(t) \\
			u_2'(t)
		\end{bmatrix}
		=
		\begin{bmatrix}
			u_2(t) \\
			f(t, u_1(t), u_2(t))
		\end{bmatrix}
	\]
	The explicit Euler update scheme discretizes time into steps $t_k = a + kh$,
	where $h = (b-a)/n$ is the uniform step size. The numerical approximation
	$\mathbf{u}_k \approx \mathbf{u}(t_k)$ advances via:
	\[
		\mathbf{u}_{k+1} = \mathbf{u}_k + h \mathbf{f}(t_k, \mathbf{u}_k)
	\]
	For the given problems, the specific governing systems become:
	\begin{enumerate}
		\item $\mathbf{u}' = [u_2, \: \sin(2t) - 9u_1]^\intercal$
		\item $\mathbf{u}' = [u_2, \: 4t + 4u_1]^\intercal$
		\item $\mathbf{u}' = [u_2, \: t - 4u_2 - 4u_1]^\intercal$
	\end{enumerate}

	\paragraph{Results}
	The integration was performed with $n = 1000$ steps for all three systems.
	The step sizes correspond to $h_1 \approx 0.00628$, $h_2 = 0.0015$, and $h_3 =
		0.004$, respectively. The results alongside their corresponding absolute error
	propagation are displayed below.

	\begin{center}
		\begin{tabular}{cc}
			\includegraphics[width=0.45\textwidth]{figures/6-2/euler-method-es3-1.pdf} &
			\includegraphics[width=0.45\textwidth]{figures/6-2/euler-method-error-es3-1.pdf} \\
			{\small Sys 1: Solution for $y'' + 9y = \sin(2t)$}                         &
			{\small Sys 1: Absolute Error Propagation}                                       \\[1.5em]
		\end{tabular}
	\end{center}

	\begin{center}
		\begin{tabular}{cc}
			\includegraphics[width=0.45\textwidth]{figures/6-2/euler-method-es3-2.pdf} &
			\includegraphics[width=0.45\textwidth]{figures/6-2/euler-method-error-es3-2.pdf} \\
			{\small Sys 2: Solution for $y'' - 4y = 4t$}                               &
			{\small Sys 2: Absolute Error Propagation}                                       \\
		\end{tabular}
	\end{center}

	\begin{center}
		\begin{tabular}{cc}
			\includegraphics[width=0.45\textwidth]{figures/6-2/euler-method-es3-3.pdf} &
			\includegraphics[width=0.45\textwidth]{figures/6-2/euler-method-error-es3-3.pdf} \\
			{\small Sys 3: Solution for $y'' + 4y' + 4y = t$}                          &
			{\small Sys 3: Absolute Error Propagation}                                       \\
		\end{tabular}
	\end{center}

	\paragraph{Discussion}
	The numerical results reflect the global truncation error characteristic of
	the first-order forward Euler method, $\mathcal{O}(h)$. Although $n=1000$
	guarantees a relatively small step size $h$, the stability of the method plays a
	critical role in the error propagation:
	\begin{itemize}
		\item \textbf{System 1 (Oscillatory):} The homogeneous part of the ODE
		      yields purely imaginary eigenvalues ($\lambda = \pm 3i$). The stability region
		      of the explicit Euler method strictly excludes the imaginary axis. Consequently,
		      the numerical solution unconditionally amplifies the oscillations over time,
		      leading to a linearly bounding envelope of the absolute error as a phase drift
		      occurs.
		\item \textbf{System 2 (Exponential Growth):} The exact solution
		      exhibits an exponentially growing term $e^{2t}$. Since the function values grow
		      rapidly, the local truncation error, which is proportional to $y''(t)$, also
		      scales exponentially. Thus, the absolute error diverges dramatically as $t \to
			      1.5$.
		\item \textbf{System 3 (Exponential Decay):} Due to the dampening factor
		      $e^{-2t}$, the exact transient solution decays toward the linear particular
		      solution. The absolute error peaks during the transient phase near $t=0.5$ and
		      subsequently decays. This behavior demonstrates that Euler's method performs
		      relatively well when the underlying system's dynamics are contractive (stable
		      eigenvalues).
	\end{itemize}

	\paragraph{Conclusion}
	The first-order Euler method captures the basic phenomenological behavior of
	the three systems for the chosen step sizes.

\end{solution}

\subsection{Runge-Kutta methods}

\begin{exercise}
	\begin{enumerate}
		\item Write a program to solve a system of differential equations using
		      the IE2 method.
		\item Write a program to solve a system of differential equations using
		      the RK4 method.
		\item Test your implementation for the IE2 and RK4 methods on the IVP
		      \[
			      u' = -2t u, \quad 0 \le t \le 2, \quad u(0) = 2; \quad \hat{u}(t) = 2e^{-t^2}
		      \]
		      Solve with $n=30,60,90,\ldots,300$ and plot the max-error
		      $||u-\hat{u}||_\infty=\max_{0\leq i\leq n} |u_i-\hat{u}(t_i)|$ as a function of
		      the \emph{number of function evaluations} in a log-log plot, together with a
		      line showing the expected rate of convergence.
		\item (\textbf{Optional}) Add the case of the ME2 and Heun's method.
	\end{enumerate}
\end{exercise}

\begin{exercise}
	In each case below, use the RK4 integrator to solve the given ODE for $0\le
		t \le 10$ with the given initial conditions. (To this end, you must find a
	sensible procedure to choose the number of steps $n$ to take.) Plot the results
	together as curves in the phase plane, that is, with $x$ and $y$ as the axes of
	the plot.
	\begin{enumerate}
		\item
		      \begin{align*}
			      x'(t) & = - 4y + x(1-x^2-y^2) \\
			      y'(t) & = 4x + y(1-x^2-y^2)
		      \end{align*}
		      with $[x(0),y(0)]=[0.1,0]$ and $[x(0),y(0)]=[0,1.9]$.
		\item
		      \begin{align*}
			      x'(t) & = - 4y - \frac{1}{4}x(1-x^2-y^2)(4-x^2-y^2) \\
			      y'(t) & = 4x - \frac{1}{4}y(1-x^2-y^2)(4-x^2-y^2)
		      \end{align*}
		      with $[x(0),y(0)]=[0.95,0]$, $[0,1.05]$, and $[-2.5,0]$.
	\end{enumerate}
\end{exercise}

\begin{exercise}
	Solve the \emph{predator-prey model} of the Lotka-Volterra equations for
	$0\leq t\leq 60$ with parameters $\alpha=0.1$, $\beta=0.25$, and initial
	conditions $y(0)=1$, $z(0)=0.01$. Plot the solutions for $y(t)$ and $z(t)$ as a
	function of $t$ on one plot, while on another figure plot the results together
	as curves in the phase plane $(y(t),z(t))$.
\end{exercise}

\subsection{Some interesting equation}

\begin{exercise}
	\begin{enumerate}
		\item Solve the Schr\"odinger equation with $V(x)=0$ considering as
		      initial wave function the Gaussian wave packet with $x_0=0$, $\sigma=1$, and
		      $k_0=0$. Consider the interval $-10\leq x\leq 10$ for the $x$ variable and
		      $0\leq t\leq 20$ for the time. Plot the absolute value squared of the wave
		      function at the times $t=0,2,3,4,5$.[\emph{Hint}: Values of $\Delta x=0.1$ and
				      $\Delta t=0.01$ are reasonable starting values for this study.]

		\item Verify that the norm of the wave function is approximately conserved in time.

		\item \textbf{(Optional)} Use the analytic solution to determine the
		      error due to the discretizations.
	\end{enumerate}
\end{exercise}

\begin{exercise}
	\begin{enumerate}
		\item Solve the Schr\"odinger equation with $V(x)=0$ considering as
		      initial wave function the Gaussian wave packet with $x_0=20$, $\sigma=2$, and
		      $k_0=2$. Consider the interval $0\leq x\leq 100$ for the $x$ variable and $0\leq
			      t\leq 20$ for the time. Plot the absolute value squared of the wave function at
		      the times $t=0,10,20$.[\emph{Hint}: Values of $\Delta x=0.1$ and $\Delta t=0.01$
				      are reasonable starting values for this study.]
		\item Verify that the norm of the wave function is approximately
		      conserved in time.
		\item \textbf{(Optional)} Use the analytic solution to determine the
		      error due to the discretizations.
	\end{enumerate}
\end{exercise}

\begin{exercise}
	\begin{enumerate}
		\item Solve the Schr\"odinger equation with the potential
		      \[
			      V(x)=
			      \begin{cases}
				      2 & \text{if}\  -1 \leq x \leq 1\,, \\[2ex]
				      0 & \text{otherwise}\,,
			      \end{cases}
		      \]
		      considering as initial wave function the Gaussian wave packet with
		      $x_0=-20$, $\sigma=2$, and $k_0=2$. Consider the interval $-100\leq x\leq 100$
		      for the $x$ variable and $0\leq t\leq 40$ for the time. Plot the absolute value
		      squared of the wave function at the times $t=0,10,30$.
		\item Compute the reflection and transmission coefficients and verify
		      that $R+T\approx 1$.[\emph{Hint}: Consider a value of $t_{\rm max}\approx 30$.]
		\item \textbf{(Optional)} Study the behavior of $R$ and $T$ as you vary
		      the height and width of the potential and/or the parameters of the wave
		      function.
	\end{enumerate}
\end{exercise}

\begin{exercise}
	Solve the time-independent Schr\"odinger equation with the potential
	\[
		V(x)=
		\begin{cases}
			-10 & \text{if}\  -1 \leq x \leq 1\,, \\[2ex]
			0   & \text{otherwise}\,.
		\end{cases}
	\]
	\begin{enumerate}
		\item Find the energies $E_n$ and wave functions $\varphi_n(x)$. How
		      many bound states do you find?[\emph{Hint}: Consider discretizing the
			      Hamiltonian in the interval $-5\leq x\leq 5$ with $\Delta x=0.05$. Use the
			      QR-algorithm for solving the eigenvalue problem with a number of iterations
			      $n\gtrsim 200$.]
		\item Plot the results for the corresponding $|\varphi_n(x_i)|^2$ and
		      binding energies together with the potential.[\emph{Hint}: Remember to normalize
				      the wave functions. Do you really need to compute the integral of the squared
				      wave function for that?]
		\item \textbf{(Optional)} Compare the results that you found in point
		      (a) for the energies $E_n$ with the exact solutions of the transcendental
		      equations for the finite square well
		      \[
			      V(x)=
			      \begin{cases}
				      -V_0 & \text{if}\  -a \leq x \leq a\,, \\[2ex]
				      0    & \text{otherwise}\,.
			      \end{cases}
		      \]
		      We recall that, given,
		      \[
			      \rho=\sqrt{-{2mE\over \hbar}}\,,
			      \qquad
			      k = \sqrt{2m(E+V_0)\over \hbar^2}\,,
			      \qquad
			      -V_0<E<0\,,
		      \]
		      the bound state solutions must satisfy,
		      \[
			      \text{Even parity:} \quad  \tan (ka) = {\rho\over k}\,;
			      \qquad
			      \text{Odd parity:} \quad \cot (ka) = -{\rho\over k}\,.
		      \]
		      Use your favorite rootfinding algorithm to solve these equations for $E$.
	\end{enumerate}
\end{exercise}
